% !TeX root = ../main.tex
教示時と同じ指のSensor\_off条件において成功率の低下およびばらつきが生じた主要因として，視覚情報のみに依存した高さ推定の限界が挙げられる．本実験環境では，側面カメラは障害物によるオクルージョンが発生し，上面カメラは深度情報を持たないRGBカメラである．そのため，積層されたキムタオルの高さ変化を画像特徴の変化としてのみ捉える必要があるが，照明条件の変動やわずかな位置ずれにより，その特徴量は不安定となる．成功しているエピソードに関しては，上面カメラが対象物の影等を捉えて画像特徴による高さ推定が成功したか，指の柔軟性が推定誤差を吸収した結果であると推察される．\par
教示時より短い指のSensor\_on条件において1枚の把持が失敗した原因は，教示データに含まれる関節角度の分布にあると考えられる．教示データには，指が短くなった分だけアームを深く下ろすような動作データは含まれていない．本手法で用いたACTアーキテクチャは，触覚情報を利用することで接触時の微調整は可能であるが，教示された関節角度の分布から大きく逸脱するような大域的な軌道修正（教示時よりもさらに深くアームを下げる動作）までは生成できない．したがって，1枚の条件での失敗は，学習ベースの手法における汎化性能の限界を示していると言える．\par
教示時より短い指のSensor\_off条件において30枚以上の高さで高い成功率を示したのは，教示データと同一の指長さでは画像特徴量による高さ推定のずれで深く掴みに行き過ぎてつかむことが難しくなっていたが，指が短くなったことでそのずれを吸収し適切な深さになったためと考えられる．